\documentclass[11pt]{article}
\usepackage[margin=1in]{geometry}
\usepackage{amsmath,amssymb,amsthm}
\usepackage{hyperref}
\usepackage{enumitem}

\newtheorem{theorem}{Theorem}
\newtheorem{lemma}[theorem]{Lemma}
\newtheorem{corollary}[theorem]{Corollary}
\theoremstyle{definition}
\newtheorem{definition}[theorem]{Definition}
\theoremstyle{remark}
\newtheorem{remark}[theorem]{Remark}

\usepackage{xurl}
\let\oldus\_
\renewcommand{\_}{\oldus\allowbreak}
\newcommand{\dl}{\operatorname{dl}}

\title{A disproof of Conjecture 00000002307\\[2pt]
\large Fitting height versus derived length}
\date{}

\begin{document}
\sloppy
\maketitle

\section{The conjecture}

Conjecture 00000002307 reads:

\begin{quote}
\emph{Definition:} The Fitting height and derived length are two invariants of solvable groups.
\emph{Conjecture:} Fitting height $\le 2\cdot(\text{derived length}) - 1$ (a linear bound); the bound is
controlled by layerwise projections of supersolvable towers; tightness is attained by extremal examples
with total derived chain $2$.
\end{quote}

It is the conjunction of three clauses. Writing $h(G)$ for the Fitting height and $\dl(G)$ for the derived
length:
\begin{enumerate}[label=(\roman*)]
  \item (\emph{linear bound}) $h(G) \le 2\,\dl(G) - 1$ for every finite solvable group $G$;
  \item the bound is ``controlled by layerwise projections of supersolvable towers'';
  \item (\emph{tightness}) the bound is attained by an example of derived length $2$, i.e.\ some solvable group
        $G$ has $\dl(G)=2$ and $h(G) = 2\,\dl(G)-1 = 3$.
\end{enumerate}
Clause (ii) is not a precise mathematical statement, and we do not need it: the conjunction
(i)$\wedge$(ii)$\wedge$(iii) is false as soon as (i) or (iii) is false. We show that \textbf{both (i) and
(iii) are false}. Clause (iii) fails for a structural reason (Theorem~\ref{thm:main}: $h(G)\le \dl(G)$ for
every solvable group), so the refutation does not depend on how the trivial group is treated.

\section{Definitions}

We use the standard definitions, as stated for instance on the Wikipedia pages ``Fitting length''
(\url{https://en.wikipedia.org/wiki/Fitting_length}) and ``Solvable group''
(\url{https://en.wikipedia.org/wiki/Solvable_group}):

\begin{definition}
A \emph{Fitting chain} (Fitting series, nilpotent series) of length $n$ of a group $G$ is a subnormal series
\[
  1 = H_0 \trianglelefteq H_1 \trianglelefteq \cdots \trianglelefteq H_n = G
\]
whose successive quotients $H_{i+1}/H_i$ are nilpotent. The \emph{Fitting length} (Fitting height,
nilpotent length) $h(G)$ is the smallest length of a Fitting chain.
\end{definition}

``A Fitting chain \dots{} for a group is a subnormal series with nilpotent quotients \dots{} The Fitting
length or nilpotent length of a group is defined to be the smallest possible length of a Fitting chain''
(Wikipedia, \emph{Fitting length}).

\begin{definition}
The derived series is $G^{(0)} = G$, $G^{(k+1)} = [G^{(k)}, G^{(k)}]$. The group is solvable if
$G^{(n)} = 1$ for some $n$, and the \emph{derived length} $\dl(G)$ is the least such $n$.
\end{definition}

``The least $n$ such that $G^{(n)} = 1$ is called the derived length of the solvable group $G$''
(Wikipedia, \emph{Solvable group}). With these definitions the trivial group has $h(1) = \dl(1) = 0$:
the chain $1 = H_0 = G$ has length $0$, and $G^{(0)} = G = 1$.

\section{Fitting height is at most derived length}

\begin{lemma}\label{lem:abelian}
Let $B \le A$ be subgroups of a group $G$ with $B \trianglelefteq A$ and $[A,A] \le B$. Then $A/B$ is abelian,
hence nilpotent.
\end{lemma}

\begin{proof}
For $x,y\in A$ we have $(xy)^{-1}(yx) = y^{-1}x^{-1}yx = [y^{-1},x^{-1}] \in [A,A] \le B$, so $xyB = yxB$.
Abelian groups are nilpotent (of class $\le 1$).
\end{proof}

\begin{theorem}\label{thm:main}
For every solvable group $G$ (finite or not), $h(G) \le \dl(G)$.
\end{theorem}

\begin{proof}
Let $d = \dl(G)$, so $G^{(d)} = 1$. Put $H_i = G^{(d-i)}$ for $0\le i\le d$. Then $H_0 = G^{(d)} = 1$ and
$H_d = G^{(0)} = G$. Each $G^{(k)}$ is normal in $G$ (it is characteristic), and
$G^{(k+1)} = [G^{(k)},G^{(k)}] \le G^{(k)}$; hence $H_i \le H_{i+1}$ and $H_i \trianglelefteq H_{i+1}$.
Writing $A = H_{i+1} = G^{(k)}$ and $B = H_i = G^{(k+1)}$ with $k = d-i-1$, we have $[A,A] = B$, so
$H_{i+1}/H_i$ is nilpotent by Lemma~\ref{lem:abelian}. Thus $(H_i)_{0\le i\le d}$ is a Fitting chain of
length $d$, and $h(G) \le d$.
\end{proof}

\begin{corollary}\label{cor}
Let $G$ be a solvable group.
\begin{enumerate}[label=(\alph*)]
  \item If $G \ne 1$ then $\dl(G)\ge 1$, hence $h(G) \le \dl(G) \le 2\,\dl(G) - 1$: clause (i) holds for
        every nontrivial solvable group.
  \item If $h(G) = 2\,\dl(G) - 1$ then $\dl(G) \le 1$. In particular no solvable group of derived length $2$
        has $h(G) = 3$.
\end{enumerate}
\end{corollary}

\begin{proof}
(a) If $\dl(G) = 0$ then $G = G^{(0)} = 1$. (b) $2\,\dl(G) - 1 = h(G) \le \dl(G)$ gives $\dl(G)\le 1$.
If $\dl(G) = 2$ then $h(G) \le 2 < 3$.
\end{proof}

\section{The refutation}

\begin{theorem}\label{thm:refute}
\begin{enumerate}[label=(\arabic*)]
  \item Clause (i) is false: the trivial group is finite and solvable with $h = \dl = 0$, and
        $0 > 2\cdot 0 - 1$.
  \item Clause (iii) is false: by Corollary~\ref{cor}(b), no solvable group (finite or infinite) of derived
        length $2$ attains the bound, which would require Fitting height $3$.
\end{enumerate}
Consequently Conjecture 00000002307 is false.
\end{theorem}

\begin{remark}
Part (2) is the substantive refutation. Even if clause (i) is read only for nontrivial groups, where it is
true (Corollary~\ref{cor}(a)), the claimed extremal examples ``with total derived chain $2$'' cannot exist:
equality $h = 2\,\dl - 1$ occurs exactly for the nontrivial abelian groups ($h = \dl = 1$). For example $S_3$
has $\dl = 2$ and $h = 2$, and $S_4$ has $\dl = 3$ and $h = 3$; in both cases $h \le \dl$, and the bound
$2\,\dl - 1$ (equal to $3$ and $5$) is not attained.
\end{remark}

\section{Lean 4 formalization}

The directory \texttt{lean/} contains a Lean~4 project (Lean \texttt{v4.33.1}, Mathlib \texttt{v4.33.1}).
The single source file \texttt{Conjecture2307/Basic.lean} contains:

\begin{itemize}
  \item \texttt{FittingChain G n}: subgroups \texttt{H 0 = $\bot$ $\le$ $\cdots$ $\le$ H n = $\top$}, each
        normal in the next (\texttt{(H i).subgroupOf (H (i+1))} is normal), with
        \texttt{Group.IsNilpotent} of the quotient group $H_{i+1}/H_i$ (\texttt{H (i+1) / (H i).subgroupOf (H (i+1))});
  \item \texttt{fittingLength G := sInf \{n | Nonempty (FittingChain G n)\}} and
        \texttt{derivedLength G := sInf \{n | derivedSeries G n = $\bot$\}} (Mathlib's
        \texttt{derivedSeries});
  \item \texttt{LinearBound} (clause (i), over all finite solvable \texttt{G : Type}, inequality in
        $\mathbb Z$) and \texttt{TightAtDerivedLengthTwo} (clause (iii));
  \item \texttt{fittingLength\_le\_derivedLength} (Theorem~\ref{thm:main}),
        \texttt{linearBound\_of\_nontrivial} and \texttt{derivedLength\_le\_one\_of\_tight}
        (Corollary~\ref{cor}),
        \texttt{not\_linearBound}, \texttt{not\_tightAtDerivedLengthTwo} (Theorem~\ref{thm:refute}), and
  \item the main theorem
\[
  \texttt{conjecture\_00000002307\_false} :\ \neg\,(\texttt{LinearBound} \wedge \texttt{TightAtDerivedLengthTwo}).
\]
\end{itemize}

There is no \texttt{sorry}, no \texttt{native\_decide} and no added axiom;
\texttt{\#print axioms conjecture\_00000002307\_false} reports only \texttt{propext},
\texttt{Classical.choice} and \texttt{Quot.sound}. To check:
\begin{verbatim}
cd lean
lake exe cache get
lake build
\end{verbatim}

\end{document}
